Because a financial peer group determines whose returns enter a firm's
peer average in a spatial model of equity returns, the distance used to
select peers shapes what the model measures. We ask when a frozen tabular
foundation model selects better peers than a distance computed directly
from the same financial records, measuring fit by the in-sample
spatial-lag log likelihood of US equity returns from 2011 to 2025.
Although TabPFN peers chosen from four financial ratios fit better than
peers chosen by Fern\'andez's characteristic-history distance, by
\ffmatched{factor-gain} nats per issuer-quarter after five Fama--French
factors and momentum, a ten-neighbour graph on the same standardized
ratios fits at least as well as TabPFN under market controls. Whereas
encoding adds little to four ratios, it adds fit once the inputs extend to
broad filing data, where all four encoders we study have positive estimated
gains over the direct graph built from the same fields and the TabPFN gain
of \ffmatched{broad-gain} nats per issuer-quarter under market controls
remains positive under multifactor controls. A four-firm example shows why
prediction accuracy cannot settle peer selection, since an invertible
recoding can preserve every prediction while changing nearest neighbours.
While the leading coordinate of the broad-input representation orders
firms by valuation, peers chosen along it have a mean Jaccard similarity of
only \ffmatched{overlap} with the full representation's peers, and
restricting the graph to that coordinate lowers fit by more than the
encoder's gain over the direct graph.
